\paragraph{40 状态真实输入核的单源临界带、Lucas 素长估值与有限部双机制}
在定理 \ref{thm:conclusion-realinput40-atomic-surgery-three-level-invertibility} 已把单原子 surgery 的 \(\zeta\)/ghost/primitive 三层反演固定、而定理 \ref{thm:conclusion-realinput40-abel-gap-positive-primitive-mass} 已把原子差额的 Abel 质量闭式锁定之后，还可把真实输入 \(40\) 状态核中由 primitive 轨道谱拆分、Dirichlet 完成截面、Lucas 型素长闭式与 Perron 有限部接口共同诱导的解析几何进一步压缩为下列五条后果。以下固定
\[
\varphi=\frac{1+\sqrt5}{2},
\qquad
F:=\begin{pmatrix}1&1\\[2pt]1&0\end{pmatrix},
\qquad
\lambda:=\rho(M)=\varphi^2,
\qquad
z_\star:=\lambda^{-1}=\varphi^{-2},
\]
并记
\[
\widehat{\zeta}_M(s):=\zeta_M(\lambda^{-s}).
\]

\begin{theorem}[开临界带的单源分解]\label{thm:conclusion-realinput40-open-strip-single-source}
\leanverified{paper\_conclusion\_realinput40\_open\_strip\_single\_source}
记
\[
H(s):=
\frac{\zeta_{F\otimes F}(\lambda^{-s})}
{(1-\lambda^{-s})^2(1+\lambda^{-s})(1-\varphi^{-1}\lambda^{-s})}.
\]
则 \(H\) 在开条带 \(0<\Re(s)<1\) 内全纯且无零，并且
\[
\widehat{\zeta}_M(s)=\frac{H(s)}{1+\varphi^{1-2s}}.
\]
特别地，\(\widehat{\zeta}_M\) 在开临界带中的全部极点都由单一因子 \(1+\varphi^{1-2s}\) 产生。
\end{theorem}
\begin{proof}
由命题 \ref{prop:real-input-40-fib-tensor}，
\[
\zeta_M(z)=\frac{1}{(1-z)^2(1+z)}\,
\zeta_{F\otimes F}(z)\,\zeta_{-F}(z),
\qquad
\zeta_{-F}(z)=\frac{1}{1+z-z^2}.
\]
又
\[
1+z-z^2=(1-\varphi^{-1}z)(1+\varphi z),
\]
故代入 \(z=\lambda^{-s}\) 得
\[
\widehat{\zeta}_M(s)
=
\frac{\zeta_{F\otimes F}(\lambda^{-s})}
{(1-\lambda^{-s})^2(1+\lambda^{-s})(1-\varphi^{-1}\lambda^{-s})(1+\varphi\lambda^{-s})}
=
\frac{H(s)}{1+\varphi^{1-2s}}.
\]
其中 \(1+\varphi\lambda^{-s}=1+\varphi^{1-2s}\)。另一方面，
\[
\zeta_{F\otimes F}(z)=\frac{1}{(1-\lambda z)(1+z)^2(1-\lambda^{-1}z)},
\]
故 \(H\) 的可能奇点只能落在 \(\Re(s)\in\{1,0,-\tfrac12,-1\}\) 上，因而在 \(0<\Re(s)<1\) 内全纯。又 \(\zeta_{F\otimes F}\) 与分母各因子都不在该开条带内取零，故 \(H\) 在该条带内无零。证毕。
\end{proof}

\begin{theorem}[临界带极点晶格的等距同留数性]\label{thm:conclusion-realinput40-open-strip-pole-lattice-equal-residue}
设
\[
s_k:=\frac12+\frac{(2k+1)\pi i}{\log\lambda},
\qquad
k\in\ZZ.
\]
则 \(\widehat{\zeta}_M\) 在开临界带中的全部极点恰为 \(\{s_k\}_{k\in\ZZ}\)；这些极点全为单极点，并满足
\[
\operatorname*{Res}_{s=s_k}\widehat{\zeta}_M(s)
=
\frac{H(s_0)}{\log\lambda}
\qquad
(k\in\ZZ).
\]
\end{theorem}
\begin{proof}
由定理 \ref{thm:conclusion-realinput40-open-strip-single-source}，开临界带内的极点恰由
\[
1+\varphi^{1-2s}=0
\]
决定。该方程等价于
\[
(1-2s)\log\varphi=(2k+1)\pi i,
\]
亦即
\[
s=s_k=\frac12+\frac{(2k+1)\pi i}{\log\lambda}.
\]
再对分母求导：
\[
\frac{\mathrm d}{\mathrm ds}\bigl(1+\varphi^{1-2s}\bigr)
=-\log\lambda\cdot \varphi^{1-2s}.
\]
在 \(s=s_k\) 处有 \(\varphi^{1-2s_k}=-1\)，故该导数等于 \(\log\lambda\neq 0\)，从而所有极点均为单极点。最后，\(H\) 只依赖于 \(\lambda^{-s}\)，因而具有虚周期 \(2\pi i/\log\lambda\)；而
\[
s_{k+1}-s_k=\frac{2\pi i}{\log\lambda},
\]
所以 \(H(s_k)=H(s_0)\) 对一切 \(k\) 成立。留数公式遂化为
\[
\operatorname*{Res}_{s=s_k}\widehat{\zeta}_M(s)
=
\frac{H(s_k)}{\log\lambda}
=
\frac{H(s_0)}{\log\lambda}.
\]
证毕。
\end{proof}

\begin{theorem}[primitive 平方根异常的奇支撑]\label{thm:conclusion-realinput40-primitive-sqrt-anomaly-odd-support}
记 \(p_n:=p_n(M)\)。则有统一渐近
\[
p_n
=
\frac{\lambda^n}{n}
-\mathbf 1_{\{n\ \mathrm{odd}\}}\frac{\lambda^{n/2}}{n}
+O\!\left(\frac{\lambda^{n/3}}{n}\right).
\]
等价地，
\[
p_n
=
\frac{\lambda^n}{n}
+\frac{\bigl((-1)^n-\mathbf 1_{\{2\mid n\}}\bigr)\lambda^{n/2}}{n}
+O\!\left(\frac{\lambda^{n/3}}{n}\right).
\]
特别地，\(\lambda^{n/2}\) 级别的平方根异常完全支撑在奇长度上；在偶长度上，M\"obius 反演中的 \(d=1\) 与 \(d=2\) 两项于该尺度发生精确抵消。
\end{theorem}
\begin{proof}
命题 \ref{prop:finite-rh-40} 已给出：真实输入 \(40\) 状态核在 \(|\mu|=\sqrt{\lambda}\) 上恰有唯一非 Perron 模态 \(-\sqrt{\lambda}\)，而其余特征值严格落在 \(|\mu|<\sqrt{\lambda}\) 内。因此定理 \ref{thm:primitive-odd-even-sqrt-cancellation} 可直接应用，并在 \(\rho=\lambda\) 时给出
\[
p_n
=
\frac{\lambda^n}{n}
-\mathbf 1_{\{n\ \mathrm{odd}\}}\frac{\lambda^{n/2}}{n}
+O\!\left(\frac{\lambda^{n/3}}{n}\right).
\]
第二个写法只是把奇偶指示函数作代数重写。最后一句正是同一定理中关于 \(d=1\) 与 \(d=2\) 精确消项机制的直接重述。证毕。
\end{proof}

\begin{theorem}[素长 primitive 层的 Lucas--Wieferich 估值律]\label{thm:conclusion-realinput40-prime-length-lucas-wieferich-valuation}
设 \(p\) 为奇素数，记
\[
L_p:=\Tr(F^p),
\qquad
p_p:=p_p(M),
\qquad
\nu_p:=v_p(L_p-1).
\]
则有精确公式
\[
p_p=\frac{L_p(L_p-1)}{p},
\qquad
L_p\equiv 1 \pmod p,
\qquad
v_p(p_p)=\nu_p-1.
\]
因而对任意 \(k\ge 0\)，
\[
p^k\mid p_p
\iff
p^{k+1}\mid (L_p-1),
\qquad
p_p\equiv \frac{L_p-1}{p}\pmod p.
\]
\end{theorem}
\begin{proof}
推论 \ref{cor:real-input-40-prime-length-product} 已给出
\[
p_p=\frac{L_p(L_p-1)}{p}.
\]
另一方面，引理 \ref{lem:pom-lucas-modp} 表明 \(L_p\equiv 1\pmod p\)，故 \(v_p(L_p)=0\)。于是
\[
v_p(p_p)=v_p(L_p)+v_p(L_p-1)-1=\nu_p-1.
\]
这立即推出
\[
p^k\mid p_p
\iff
v_p(p_p)\ge k
\iff
\nu_p\ge k+1
\iff
p^{k+1}\mid(L_p-1).
\]
又因 \(L_p\equiv 1\pmod p\)，故在模 \(p\) 意义下 \(L_p\) 为单位并且同余于 \(1\)，从而
\[
p_p=\frac{L_p(L_p-1)}{p}\equiv \frac{L_p-1}{p}\pmod p.
\]
证毕。
\end{proof}

\begin{theorem}[Perron 极留数的局域手术律与 Abel 有限部的三因子加法律]\label{thm:conclusion-realinput40-perron-residue-local-surgery-finitepart-additivity}
将 \(\zeta_M\) 写为
\[
\zeta_M(z)=\zeta_{F\otimes F}(z)\,R_{\mathrm{reg}}(z),
\qquad
R_{\mathrm{reg}}(z):=\frac{\zeta_{-F}(z)}{(1-z)^2(1+z)}
=\frac{1}{(1+z-z^2)(1-z)^2(1+z)}.
\]
则有：
\begin{enumerate}
  \item \(R_{\mathrm{reg}}\) 在 \(z=z_\star\) 处解析且非零，并满足显式闭式
  \[
  R_{\mathrm{reg}}(z_\star)=\frac{15+7\sqrt5}{20}.
  \]
  \item 若记
  \[
  C_M:=\lim_{z\to z_\star}(1-\lambda z)\zeta_M(z),
  \qquad
  C_{F\otimes F}:=\lim_{z\to z_\star}(1-\lambda z)\zeta_{F\otimes F}(z),
  \]
  则
  \[
  C_M=R_{\mathrm{reg}}(z_\star)\,C_{F\otimes F}.
  \]
  \item 若按注 \ref{rem:real-input-40-logM-artin-factorization} 记 \(\log\mathfrak{M}_{\mathrm{even}}\) 为 \(\zeta_{\mathrm{even}}(z)=(1-3z+z^2)^{-1}\) 的有限部常数，并定义
  \[
  \log\mathfrak{M}_{\mathrm{twist}}:=K_{-F},
  \qquad
  \log\mathfrak{M}_{\mathrm{triv}}:=P_{\mathrm{triv}}(z_\star),
  \]
  则 Abel 有限部满足
  \[
  \log\mathfrak{M}
  =
  \log\mathfrak{M}_{\mathrm{even}}
  +\log\mathfrak{M}_{\mathrm{twist}}
  +\log\mathfrak{M}_{\mathrm{triv}}.
  \]
\end{enumerate}
因此，Perron 主极点的留数修正满足纯局域的乘法手术律，而 Abel 有限部则保持 even/twist/triv 三因子的全局加法律。
\end{theorem}
\begin{proof}
命题 \ref{prop:real-input-40-fib-tensor} 已给出
\[
\zeta_M(z)=\frac{1}{(1-z)^2(1+z)}\,\zeta_{F\otimes F}(z)\,\zeta_{-F}(z),
\]
故所示 \(R_{\mathrm{reg}}\) 分解成立。由于 \(z_\star=\varphi^{-2}\in(0,1)\) 且 \(\zeta_{-F}\) 的极点位于 \(z=-\varphi^{-1}\) 与 \(z=\varphi\)，故 \(R_{\mathrm{reg}}\) 在 \(z_\star\) 处解析且非零。直接代入并化简得
\[
R_{\mathrm{reg}}(z_\star)
=
\frac{1}{(1+\varphi^{-2}-\varphi^{-4})(1-\varphi^{-2})^2(1+\varphi^{-2})}
=
\frac{15+7\sqrt5}{20}.
\]
于是第二条只需把 \(\zeta_M=\zeta_{F\otimes F}R_{\mathrm{reg}}\) 乘以 \(1-\lambda z\) 后令 \(z\to z_\star\) 即得：
\[
C_M
=
\lim_{z\to z_\star}(1-\lambda z)\zeta_{F\otimes F}(z)\,R_{\mathrm{reg}}(z)
=
R_{\mathrm{reg}}(z_\star)\,C_{F\otimes F}.
\]
最后，注 \ref{rem:real-input-40-logM-artin-factorization} 已明确给出
\[
\log\mathfrak{M}
=
\log\mathfrak{M}_{\mathrm{even}}+K_{-F}+P_{\mathrm{triv}}(z_\star).
\]
将 \(K_{-F}\) 与 \(P_{\mathrm{triv}}(z_\star)\) 重新记作 \(\log\mathfrak{M}_{\mathrm{twist}}\) 与 \(\log\mathfrak{M}_{\mathrm{triv}}\)，便得到第三条。证毕。
\end{proof}
